\section{Related Work}
\label{sec:related_work}

\paragraph{Models for Document Understanding.}
There are two main types of language models for document understanding: (1) OCR-dependent models, which use Optical Character Recognition (OCR) to extract text for processing, including the LayoutLM series~\cite{LayoutLMv1, LayoutLMv2, LayoutLMv3} and text-only LLMs~\cite{Llama3, Qwen2}; and (2) end-to-end models, which use a visual encoder to extract features from document images, integrating them with text for input into language model backbones. Most of document-related LVLMs fall into this category, such as GPT4o~\cite{GPT4o}, Gemini-1.5~\cite{Gemini15}, Claude-3.5-Sonnet~\cite{Claude35Sonnet}, mPLUG-DocOwl2~\cite{mPLUG-DocOwl2}, Qwen2-VL~\cite{Qwen2-VL}.

\paragraph{Methods for Long Document Understanding.}
To address the challenges of cross-page document understanding with excessively long context lengths, early approaches employed hierarchical encoding methods~\cite{MP-DocVQA,MPDoc-SAS,UVL-TG,MPDoc-RMT}. In these approaches, an independent encoder processes the OCR text and visual modal information for each page, which is then passed to a small language model decoder for cross-page contextual learning. However, this approach is limited by the redundancy in OCR inputs, which restricts the context length and leads to the accumulation of errors~\cite{LayoutLMv1,LayoutLMv2,DocFormer}. Recently, with the rise of multimodal large models, methods based on multimodal Retrieval-Augmented Generation (MM-RAG)~\cite{VisRAG,GRAM,MMVQA,CFRet-DVQA,CREAM,M3DocRAG} and end-to-end multi-page large models~\cite{MANTIS,LLaVA-Next-Interleave,Leopard} have emerged. These models leverage the world knowledge of large language models to enhance understanding. End-to-end approaches mitigate error accumulation by dynamically reducing the number of visual tokens across multiple pages/images and build large scale instruction turning dataset~\cite{MANTIS,LLaVA-Next-Interleave,Leopard,mPLUG-DocOwl2}, allowing for longer context lengths. Methods such as multi-page RAG facilitate dynamic interactions with OCR and other text information to remove redundant multimodal tokens.

\vspace{-0.4mm}
\paragraph{Benchmarks for Long Document Understanding.}
Multi-page or long documents place higher demands on the model's capabilities in cross-page understanding. Current multi-page document benchmarks, such as MP-DocVQA and DUDE, do not assess documents exceeding 20 pages. MMLongBench-Doc~\cite{MMLongBench-Doc} and M-Longdoc~\cite{M-Longdoc-arxiv} have been proposed to evaluate the understanding capabilities of longer documents, which have an average of 47.5 and 210.8 pages per document, respectively. Meanwhile, MMVQA~\cite{MMVQA} and MVQA~\cite{MVQA} are proposed to better evaluate the retrieval-based method. WebQuest~\cite{WebQuest} focuses on evaluating models' performance on text-rich web images.

